% Source types: Mahbubur Rahman, Chapter 5 PDF pp.43--45; Python literals.
\booksubheading{Literal: কোডে সরাসরি লেখা মান}

কোডে সরাসরি লেখা \py{13}, \py{2.5}, \py{"ICT"} প্রভৃতি মানকে
literal বলা হয়। \py{marks = 75}-এ \py{marks} একটি নাম এবং
\py{75} একটি integer literal। একটি পূর্ণসংখ্যা বিভিন্ন ভিত্তিতে লেখা গেলেও
তার data type \py{int} থাকে।

\begin{tabularx}{\linewidth}{|l|Y|Y|}
\hline
\rowcolor{softBlue}\textbf{ধরন} & \textbf{Python-এ লেখা} & \textbf{দশমিক মান}\\\hline
দশমিক & \py{13} & 13\\\hline
বাইনারি & \py{0b1101} & 13\\\hline
অক্টাল & \py{0o15} & 13\\\hline
হেক্সাডেসিমাল & \py{0xD} & 13\\\hline
Scientific notation & \py{2.5e-3} & \(2.5\times10^{-3}=0.0025\)\\\hline
\end{tabularx}

\begin{pythoncode}
print(13, 0b1101, 0o15, 0xD)
print(2.5e-3, 1.2e4)
print(type(0xD).__name__, type(1.2e4).__name__)
\end{pythoncode}
\textbf{Output:}
\begin{codebox}\begin{english}\ttfamily
13 13 13 13\\
0.0025 12000.0\\
int float
\end{english}\end{codebox}

\textbf{পরীক্ষায় মনে রাখবে:} \py{0b}, \py{0o} ও \py{0x} হলো
integer literal-এর ভিত্তি নির্দেশক prefix। বৈজ্ঞানিক রূপে \py{e} বা
\py{E}-এর পরের পূর্ণসংখ্যাটি 10-এর ঘাত বোঝায়। \py{1.2e4} লেখায়
দশমিক বিন্দু আছে এবং মানটি \py{float}। String হিসেবে লেখা
\py{"0b1101"} এবং সংখ্যামান \py{0b1101} এক ধরনের data নয়।

